\section{Comandi per l'Architettura di
Rete}\label{comandi-per-larchitettura-di-rete}

Per andare a fondo nei dettagli implementativi e spiegare meglio alcuni
concetti, sono riportati alcuni passaggi fondamentali nella
configurazione della rete dell'host che esegue BlueAgent, o di comandi
eseguiti direttamente da BlueAgent.

\subsection{Configurazione OVN e Open
vSwitch}\label{configurazione-ovn-e-open-vswitch}

Qui riporto la configurazione per implementare la SDN. Il \textbf{piano
dei dati} è implementato da Open vSwitch (\textbf{OVS}) e deve esser
configurato in modo che possa comunicare con Incus tramite l'interfaccia
di OVN

Si configura OVS in modo che possa comunicare con l'interfaccia south
bound di OVN (\passthrough{\lstinline!ovnsb\_db.sock!}):

\begin{lstlisting}[language=bash]
sudo ovs-vsctl set open_vswitch . \
    external_ids:ovn-remote=unix:/run/ovn/ovnsb_db.sock \
    external_ids:ovn-encap-type=geneve \
    external_ids:ovn-encap-ip=127.0.0.1
\end{lstlisting}

\begin{itemize}
\tightlist
\item
  \textbf{\passthrough{\lstinline!ovn-encap-type=geneve!}}:
  Specifichiamo Geneve come protocollo da utilizzare per il tunneling.
  Questo passaggio è necessario anche se non utilizziamo mai Geneve
\item
  \textbf{\passthrough{\lstinline!ovn-encap-ip=127.0.0.1!}}: Imposta la
  terminazione del tunnel nel localhost.
\end{itemize}

Si configura invece l'interfaccia North Bound in modo che possa
comunicare con Incus:

\begin{lstlisting}[language=bash]
ovn-nbctl set-connection ptcp:6641:127.0.0.1
incus config set network.ovn.northbound_connection tcp:127.0.0.1:6641
\end{lstlisting}

\begin{quote}
\textbf{OVN Northbound}: è l'interfaccia dove dichiaro le reti che
voglio creare e le proprietà che devono rispettare. \textbf{OVN
Southbound DB:} memorizza le tabelle di inoltro e la logica da applicare
sul piano dei dati.
\end{quote}

\subsection{\texorpdfstring{\(\text{production-}k\) Routing e
Connettività dei Peer (WireGuard e
nftables)}{\textbackslash text\{production-\}k Routing e Connettività dei Peer (WireGuard e nftables)}}\label{textproduction-k-routing-e-connettivituxe0-dei-peer-wireguard-e-nftables}

L'accesso dei partecipanti ai container della rete
\(\text{production-}k\) richiede l'applicazione di regole di inoltro
esplicite tra \(\text{wg-}k\) e \(\text{br-team-}k\), da applicare nel
kernel Linux mediante \passthrough{\lstinline!nftables!}.

Le regole per il traffico vengono definite dentro la tabella
\passthrough{\lstinline!inet blueagent!}:

\begin{lstlisting}[language=bash]
sudo nft add table inet blueagent
sudo nft add set inet blueagent team_links '{ type ifname . ifname; flags interval; }'

sudo nft add element inet blueagent team_links { "wg-$k" . "br-team-$k", "br-team-$k" . "wg-$k" }
\end{lstlisting}

All'interno della catena \passthrough{\lstinline!forward!}, prima di
accettare il traffico, viene forzato il \emph{clamping} del MSS per il
traffico TCP per prevenire blackhole dovuti alla riduzione di MTU dei
tunnel. Subito dopo, il traffico lecito tra membri e container dello
stesso team viene autorizzato verificando l'appartenenza al set
\passthrough{\lstinline!@team\_links!}:

\begin{lstlisting}[language=bash]
sudo nft add chain inet blueagent forward '{ type filter hook forward priority filter; policy accept; }'
sudo nft add rule inet blueagent forward tcp flags syn / syn,rst tcp option maxseg size set rt mtu
sudo nft add rule inet blueagent forward iifname . oifname @team_links accept
\end{lstlisting}

A priori, si vuole impedire l'inoltro dei pacchetti applicando l'azione
\passthrough{\lstinline!drop!} tra le interfacce
\passthrough{\lstinline!wg!} e \passthrough{\lstinline!br-team!} non
autorizzate, in modo da bloccare rotte e comunicazioni incrociate non
esplicitamente accettate nella tabella:

\begin{lstlisting}[language=bash]
sudo nft add rule inet blueagent forward iifname "wg-*" oifname "br-team-*" drop
sudo nft add rule inet blueagent forward iifname "br-team-*" oifname "wg-*" drop
sudo nft add rule inet blueagent forward iifname "br-team-*" oifname "br-team-*" drop
sudo nft add rule inet blueagent forward iifname "wg-*" oifname "wg-*" drop
sudo nft add rule inet blueagent forward iifname "br-team-*" accept
sudo nft add rule inet blueagent forward oifname "br-team-*" ct state established,related accept
\end{lstlisting}

Per consentire l'accesso a Internet ai container e ai peer, viene
applicato il mascheramento dell'indirizzo sorgente. Si deve dunque
abilitare il \passthrough{\lstinline!nat!} (il meccanismo di traduzione
degli indirizzi) poco prima che il frame esca e venga immesso
fisicamente sulla porta. Infine si specifica che si applica il nat
dinamico (\passthrough{\lstinline!masquerade!}) per gli indirizzi che
provengono dalla rete dei container verso l'esterno.

\begin{lstlisting}[language=bash]
sudo nft add chain inet blueagent postrouting '{ type nat hook postrouting priority srcnat; policy accept; }'
sudo nft add rule inet blueagent postrouting ip saddr 10.0.0.0/8 ip daddr != 10.0.0.0/8 masquerade
\end{lstlisting}
